\chapter{Finite-Rank Gauge Constructions and Locality Controls}
\label{ch:finite-rank-gauge-locality}

\section{Why exact control families matter}

The preceding chapters isolate the representation and alignment operations that enter an
operator comparison.  This chapter assembles a second proof arc needed by higher-rank
applications: an exact finite-rank parent with a qualified retained sector, an arbitrary
right-unitary gauge construction, and finite Fourier statements that distinguish spectrum
preservation from representation locality.  The objects are source-neutral controls.  They do
not identify a runtime DFT or Wannier subspace, certify floating-point arrays, or establish the
scientific adequacy of a material model.

Throughout this chapter, ``rank $N$'' means that the retained coordinate space has dimension
$N$.  It does not mean that a represented matrix has algebraic rank $N$ or is nonsingular.

\section{A qualified finite-rank analytic parent}

Let $r\geq 2$ be the retained dimension.  For an index $j$ with $0\leq j<r$, define the evenly
spaced orbital coordinate
\[
  x_j=\frac{2j-r+1}{r-1}.
\]
Then $-1\leq x_j\leq 1$, and adjacent coordinates satisfy
\[
  x_{j+1}-x_j=\frac{2}{r-1}.
\]
For a nonempty finite bond inventory $B$ with phases $\phi_b$, define
\[
  h(\phi)=\frac34\cos\phi+\frac14\cos(2\phi),
  \qquad
  s=\frac{1}{|B|}\sum_{b\in B}h(\phi_b).
\]
Each harmonic and therefore their average lies in $[-1,1]$.  Given a parent scale
$0\leq p\leq 1$, the retained and exterior energies are
\[
  E_{\mathrm{ret}}(x,p,s)
    =\frac25x+\frac14p\left(1+\frac15x\right)s
\]
and
\[
  E_{\mathrm{ext}}(x,p,s)
    =3+\frac3{10}x+\frac1{10}ps.
\]

\begin{theorem}[Exact finite-rank gap controls]
\label{thm:finite-rank-gap-controls}
For $x,y,s\in[-1,1]$ and $p\in[0,1]$,
\[
  E_{\mathrm{ret}}(x,p,s)\leq\frac7{10},
  \qquad
  E_{\mathrm{ext}}(y,p,s)\geq\frac{13}{5},
\]
and hence
\[
  E_{\mathrm{ext}}(y,p,s)-E_{\mathrm{ret}}(x,p,s)
    \geq\frac{19}{10}.
\]
If $r\geq2$ and $j+1<r$, adjacent selected retained energies obey
\[
  E_{\mathrm{ret}}(x_{j+1},p,s)-E_{\mathrm{ret}}(x_j,p,s)
    \geq \frac{7}{10(r-1)}>0.
\]
The universal energy and external-gap bounds are checked by
\decl{BG-FOUNDATION-THM-029--031}; the generic adjacent-gap bound is
\decl{BG-FOUNDATION-THM-034}; and the finite-bond external-gap and adjacent-gap
specializations are \decl{BG-FOUNDATION-THM-035--037}.
\end{theorem}

\begin{proof}
Since $p\in[0,1]$ and $s\in[-1,1]$, their product $q=ps$ lies in $[-1,1]$.
The retained energy can be rearranged as
\[
  E_{\mathrm{ret}}(x,p,s)
   =\left(\frac25+\frac1{20}q\right)x+\frac14q.
\]
The coefficient of $x$ is nonnegative.  Its maximum over $x,q\in[-1,1]$ is therefore
attained at $x=q=1$, giving $7/10$.  The exterior energy is affine in both $y$ and
$q$; its minimum on the same rectangle is attained at $y=q=-1$, giving $13/5$.
Subtracting the two bounds gives $19/10$.

For adjacent retained levels, direct subtraction factors as
\[
  E_{\mathrm{ret}}(x_{j+1},p,s)-E_{\mathrm{ret}}(x_j,p,s)
  =(x_{j+1}-x_j)
    \left(\frac25+\frac1{20}q\right).
\]
The first factor is $2/(r-1)$ and the second is at least $7/20$, which gives the
stated lower bound.  The Lean proof separates the interval bounds, coordinate identity,
and affine factorization so that no numerical sampling or trigonometric approximation enters the
argument.
\end{proof}

\begin{remark}
The parent is a constant exact control family.  The theorem does not prove that a consuming
Hamiltonian has these energies, that an external mesh realizes the stated structure factor, or
that an eigensolver has ordered its output consistently.
\end{remark}

\section{Exact arbitrary-dimensional right-unitary gauges}

Let $F$ be an ambient-by-retained complex frame with $F^\dagger F=I$, let $D$ be an arbitrary
retained operator, and let $G$ be a retained-space unitary.  Define
\[
  F'=FG,
  \qquad
  H'=G^\dagger D G.
\]
This fixes the action side: gauges act on frame columns from the right, and represented operators
transform by right-gauge conjugation.

\begin{theorem}[Exact finite-rank gauge construction]
\label{thm:exact-finite-rank-gauge-construction}
The constructed frame and operator satisfy
\[
  F'^\dagger F'=I,
  \qquad
  \operatorname{spec}_{\mathbb C}(H')
    =\operatorname{spec}_{\mathbb C}(D),
\]
and the known inverse $G^\dagger$ recovers both reference objects:
\[
  F'G^\dagger=F,
  \qquad
  G H'G^\dagger=D.
\]
These conclusions are checked by \decl{BG-ALIGNMENT-THM-069} and
\decl{BG-ALIGNMENT-THM-071}.
\end{theorem}

\begin{proof}
Orthonormality follows from
\[
  F'^\dagger F'=G^\dagger F^\dagger F G=G^\dagger G=I.
\]
The spectrum equality is unitary similarity over $\mathbb C$; it is an equality of spectra as
sets, not an assertion about an external eigenvalue ordering.  Finally,
\[
  F'G^\dagger=FGG^\dagger=F
\]
and
\[
  GH'G^\dagger=GG^\dagger DGG^\dagger=D.
\]
The formal theorem packages the four conclusions with the exact supplied unitary and frame
orthonormality as hypotheses.
\end{proof}

This theorem supplies an arbitrary-retained-dimension construction, while the earlier rank-two
examples expose particular locality effects.  It does not claim that every coordinate is mixed,
that the gauge was recovered by a numerical Procrustes solve, or that two externally produced
frames span the same subspace.

\section{Pointwise gauges as finite Fourier convolution}

Let $K$ be a finite abelian momentum group and $\widehat K$ its character group.  The convention
used here is
\[
  \widehat A(\chi)
    =\frac1{|K|}\sum_{k\in K}\overline{\chi(k)}A(k),
  \qquad
  A(k)=\sum_{\chi\in\widehat K}\chi(k)\widehat A(\chi).
\]
For a pointwise unitary gauge $U(k)$, set
\[
  A^U(k)=U(k)^\dagger A(k)U(k).
\]
Complex conjugation of a coefficient negates its character mode.  Applying the finite product
convolution law twice therefore gives the exact triple convolution
\[
  \widehat{A^U}(m)
   =\sum_{-a+b+c=m}
      \widehat U(a)^\dagger\,\widehat A(b)\,\widehat U(c).
\]
This is \decl{BG-FOURIER-GAUGE-THM-001}.

If $\widehat U$ is supported on $S_U$ and $\widehat A$ on $S_A$, cancellation may make the actual
support smaller, but every possible nonzero mode lies in
\[
  \operatorname{supp}\widehat{A^U}
  \subseteq (-S_U+S_A)+S_U.
\]
The formal result \decl{BG-FOURIER-GAUGE-THM-002} is deliberately an inclusion rather than an
equality.

Suppose now that $\rho:\widehat K\to\mathbb R_{\geq0}$ is invariant under inversion and
subadditive:
\[
  \rho(-m)=\rho(m),
  \qquad
  \rho(m+n)\leq\rho(m)+\rho(n).
\]
If the gauge coefficients lie in the complete shell of radius $R_U$ and the field coefficients
in the complete shell of radius $R_A$, then every mode in the triple support satisfies
\[
  \rho(-a+b+c)\leq R_A+2R_U.
\]
Consequently
\[
  \operatorname{supp}\widehat{A^U}
  \subseteq \operatorname{Shell}(R_A+2R_U),
\]
which is checked by \decl{BG-FOURIER-GAUGE-THM-003}.  This is a finite support bound.  It is not a
decay estimate and need not be sharp.

\section{Source-neutral lattice radii}

The square, triangular, simple-cubic, and normalized FCC integer shell indices admit exact
Euclidean realizations.  Their associated nonnegative radii are norms of additive real embeddings.
The formal development proves that the real square of each radius is exactly the corresponding
integer shell index (\decl{BG-FOURIER-SUPPORT-THM-001--004}), that each radius is inversion
invariant (\decl{BG-FOURIER-SUPPORT-LEM-003--006}), and that each satisfies the triangle
inequality (\decl{BG-FOURIER-SUPPORT-THM-005--008}).  Thus the general support argument gives the complete
shell inclusions
\[
  (-\operatorname{Shell}(R_U)+\operatorname{Shell}(R_A))
       +\operatorname{Shell}(R_U)
  \subseteq \operatorname{Shell}(R_A+2R_U)
\]
for all four source-neutral integer lattices
(\decl{BG-FOURIER-SUPPORT-THM-013--016}).

These are infinite integer-lattice statements.  They do not identify a finite cyclic character
with a preferred centered representative across a quotient seam.  Finite quotient radii, FFT
ordering, Nyquist conventions, reciprocal bases, and physical distance require separate bridges.

\section{Fourier truncation as exact finite least squares}

Let $S\subseteq\widehat K$ be a declared coefficient support and let $T_S\widehat A$ retain exactly
the coefficients in $S$.  If $C$ is any coefficient field supported on $S$, finite Parseval and a
coefficient-space Pythagorean identity give
\[
\begin{split}
 \sum_{k\in K}\left\|A(k)-\mathcal F^{-1}C(k)\right\|^2
 ={}&\sum_{k\in K}
       \left\|A(k)-\mathcal F^{-1}(T_S\widehat A)(k)\right\|^2\\
 &+|K|\sum_{\chi\in\widehat K}
       \left\|(T_S\widehat A)(\chi)-C(\chi)\right\|^2.
\end{split}
\]
The second term is nonnegative and vanishes exactly when $C=T_S\widehat A$.  Hence exact Fourier
truncation is the unique coefficient field attaining the equal-weight finite least-squares minimum
inside the declared support class.  These are \decl{BG-FOURIER-FINITE-THM-004} and
\decl{BG-FOURIER-FINITE-THM-005}; their arbitrary-rank periodic one-, two-, and three-dimensional
matrix specializations are \decl{BG-FOURIER-LATTICE-THM-001--005}.

For a matrix coefficient field $C$, total entrywise Frobenius power splits directly into the
identity coefficient and every nonidentity coefficient:
\[
  \sum_{x\in K}\sum_{i,j}|C(x)_{ij}|^2
  =\sum_{i,j}|C(0)_{ij}|^2
   +\sum_{x\neq0}\sum_{i,j}|C(x)_{ij}|^2.
\]
The nononsite term is defined by the second direct sum, not by subtracting separately rounded
runtime totals.  The exact partition and lattice specializations are
\decl{BG-FOURIER-FINITE-THM-006} and \decl{BG-FOURIER-LATTICE-THM-006--009}.

\section{Projection can destroy finite support}

Gauge-induced support growth is distinct from support growth caused by changing the retained
subspace.  The checked two-band Laurent example makes that distinction exact.  A nearest-neighbor
Laurent-polynomial parent $H(z)$ has finite support, but its positive dispersion contains a square
root that is not a Laurent polynomial when the product of the gap and hopping parameters is
nonzero.  The algebraic lower-band projector
\[
  P_-(k)=\frac12\left(I-\frac{H(k)}{d(k)}\right)
\]
is idempotent, lies in the lower eigenspace, and satisfies
\[
  P_-(k)H(k)P_-(k)=\frac12\bigl(H(k)-d(k)I\bigr).
\]
The final expression has no finite matrix Laurent representation although the parent does.
Therefore spectral projection can destroy finite hopping support
(\decl{BG-PROJECTION-THM-003--008}).

This is one exact counterexample, not a universal decay theorem.  It neither quantifies an
exponential tail nor establishes a Wannier obstruction, a topological invariant, or material
localization.

\section{Combined interpretation}

The results of this chapter establish three separate facts.  First, exact finite-rank control
families can have explicit retained and external gaps.  Second, right-unitary gauge changes
preserve frame orthonormality and complex set-spectrum while changing represented coordinates.
Third, both pointwise gauges and spectral projection can alter Fourier locality, by different
mechanisms.  Exact least squares then identifies the best approximation inside a frozen finite
support class and the direct power partition separates onsite from nononsite coefficients.

None of these implications binds an external array to the formal objects.  Runtime use still
requires authentication of the momentum mesh and ordering, retained basis and frames, action side,
Fourier normalization, matrix axes, support inventory, and exact or rigorously enclosed entries.
Changing projectors or performing disentanglement additionally requires hypotheses or error bounds
that are absent from the fixed-subspace gauge identities.
